\documentclass[11pt,a4paper]{article}
\usepackage{gtu-el-yazisi}
\begin{document}
\notbilgi{Görüntü İşleme}{Giriş, görüntü gösterimi ve PGM formatı (Ders 1)}{Ders1.pdf}

\sayfa{1}
\noindent Image restoration\\
\hspace*{4em}enhancement $\leftarrow$ ilk bir kaç hafta bunla uğraşacağız\\
\hspace*{4em}Segmentation\\
\hspace*{4.5em}Analyzis\\
\hspace*{0.5em}Feature Extraction\\
\hspace*{6em}Segmentation

\elcizimi[0.55]{sekiller/Ders1-s1-1.png}[1]

\noindent $(i,j)$ \qquad\quad for $i=0,\dots,M-1$\\
\hspace*{9em}for $j=0,\ \longrightarrow\ N-1$

\elcizimi[0.5]{sekiller/Ders1-s1-2.png}[1]

Bu derste aksi belirtmedikçe (0,255) arasında \okunamadi{varsayacağız}

\elcizimi[0.25]{sekiller/Ders1-s1-3.png}[1]
\elcizimi[0.4]{sekiller/Ders1-s1-4.png}[1]

\underline{boyutunu} bilmeliyiz görüntümüz,

\sayfa{2}
\elcizimi[0.35]{sekiller/Ders1-s2-1.png}[2]

$100 \times 100 \;\nearrow\;$ string ifade ederse

$10\,000 \times 4$ byte \quad 40 kb

byte byte saklayabiliriz \quad o zaman her piksel bir byte ile ifade edilir.

\underline{\underline{.pgm}} $\rightarrow$ ile nasıl gösteririz???

\noindent Magic number\\
\hspace*{1em}P2 $\rightarrow$ string demek\\
\hspace*{1em}P5 $\rightarrow$ binary\\
\hspace*{1em}Biz binary ile çalışacağız\\
\hspace*{1em}başlık kısmı string (header)

\noindent P2\\
\# my.pgm \quad $\nearrow$ bunu görsek bunu dikkate almayacağız\\
\fbox{24 \ 7} $\rightarrow$ boyutu\\
\fbox{15}\\
$\downarrow$ Önce yükseklik sonra genişlik \quad (Önce satır-sonra sütun)\\
$\downarrow$ 0'dan 15'e kadar\\
Ama derste sadece 255 yapacağız

\noindent image processing place com\\
irfan view

\sayfa{3}
\elcizimi[0.62]{sekiller/Ders1-s3-1.png}[3]

\elcizimi[0.37]{sekiller/Ders1-s3-2.png}[3]

\section{Connectivity (bağlılık)}

\elcizimi[0.37]{sekiller/Ders1-s3-3.png}[3]

0,1 şeklinde düşünüyoruz

\elcizimi[0.48]{sekiller/Ders1-s3-4.png}[3]

Manhattan distance

Manhattan böyle

\elcizimi[0.42]{sekiller/Ders1-s3-5.png}[3]

kitabın 3. bölümünden

\end{document}
